\Needspace{10\baselineskip}
\centerline{\textsc{Neuer Komet.}}

\medskip
Monatliche Correspondenz. Band IX. S. 432---435. Mai 1804.

\medskip
Der von Dr. \textsc{Olbers} am 12. März entdeckte neue Komet wurde auch in Marseille auf der Sternwarte der Marine am 7. März von \textsc{Pons} entdeckt. Am 10. März entdeckte ihn auch der Astronom \textsc{Bouvard} von der Pariser Nationalsternwarte. Um $15^u\ 57^m$ mittlerer Zeit ermittelte er eine Rektaszension von $220^\circ\ 4'$ und eine südliche Deklination von $1^\circ\ 41'$. Am 11. März um $15^u\ 39^m$ mittlerer Zeit fand er Rektaszension $220^\circ\ 10'$ und nördliche Deklination $2^\circ\ 44'$. Am 29. März um $9^u$ mittlerer Zeit fand er Rektaszension $218^\circ\ 40'$ und nördliche Deklination $46^\circ\ 37'$.

Auch \textsc{Messier} verfolgte diesen Kometen in Paris. Am 17. März um $11^u\ 34^m\ 18^s$ wahrer Zeit beobachtete er Rektaszension $220^\circ\ 29'\ 45''$ und nördliche Deklination $26^\circ\ 18'\ 24''$.

\Needspace{18\baselineskip}
Dr. \textsc{Olbers} beobachtete diesen Kometen an zehn Tagen bis zum 1. April. Aus diesen Beobachtungen bestimmte Dr. \textsc{Gauss} die folgenden parabolischen Elemente seiner Bahn:

\begin{center}
\begingroup\fontsize{9}{11}\selectfont\setlength{\tabcolsep}{3pt}
\begin{tabular}{lr}
Zeit des Perihels, 13. Februar 1804 . . . . . & $14^u\ 49^m\ 51^s$ mittl. Seeberger Zeit \\
Logarithmus des Abstands . . . . . . . . & $0.0298575$ \\
\end{tabular}\endgroup
\end{center}
{\fontsize{8}{10}\selectfont
\[\left.\begin{array}{ll}\text{Länge des Perihels . . . . . . . . . .}&148^\circ\ 44'\ 51''\\\text{Länge des aufsteigenden Knotens . . . . . .}&176\ 47\ 58\end{array}\right\}\begin{array}{l}\text{beide vom mittleren}\\\text{Äquinoktium aus gezählt}\end{array}\]}
\begin{center}
\begingroup\fontsize{9}{11}\selectfont\setlength{\tabcolsep}{3pt}
\begin{tabular}{lr}
Neigung der Bahn . . . . . . . . . . . & $56\ 28\ 40$ \\
Bewegung . . . . . . . . . . . . . . & rechtläufig. \\
\end{tabular}\endgroup
\end{center}

Mit diesen Elementen ergibt der Vergleich mit sämtlichen Beobachtungen von \textsc{Olbers} Folgendes:

\begin{center}
\begingroup\fontsize{9}{11}\selectfont\setlength{\tabcolsep}{3pt}\renewcommand{\arraystretch}{1.12}
\begin{tabular}{r|r|r|r}
1804 & \shortstack{Mittlere Zeit\\Bremen} & \multicolumn{2}{c}{\shortstack{Unterschied zwischen Rechnung\\und Beobachtung}} \\ \cline{3-4}
 & & in Rekt. & in Dekl. \\ \hline
März 12 & $12^u\,56^m\,13^s$ & $+9''$ & $0''$ \\
13 & $11\,40\,43$ & $-19$ & $-7$ \\
14 & $12\,22\,26$ & $+22$ & $+92$ \\
15 & $8\,54\,41$ & $-31$ & $-79$ \\
20 & $9\,22\,52$ & $-247 ::$ & $+121 ::$ \\
\end{tabular}\endgroup
\end{center}
\begin{center}
\begingroup\fontsize{9}{11}\selectfont\setlength{\tabcolsep}{3pt}\renewcommand{\arraystretch}{1.12}
\begin{tabular}{r|r|r|r}
1802 & \shortstack{Mittlere Zeit\\Bremen} & \multicolumn{2}{c}{\shortstack{Unterschied zwischen Rechnung\\und Beobachtung}} \\ \cline{3-4}
 & & in Rekt. & in Dekl. \\ \hline
März 22 & $8^u\,59^m\,13^s$ & $-143''$ & $+85''$ \\
27 & $8\,59\,43$ & $-55$ & $-152$ \\
28 & $8\,28\,2$ & $-27$ & $+32$ \\
29 & $8\,45\,41$ & $-26$ & $+33$ \\
April 1 & $9\,1\,52$ & $+28$ & $-1$ \\
\end{tabular}\endgroup
\end{center}

\noindent Dabei verlaufen die Fehler besonders in der Deklination so unregelmäßig, dass man sie mit gutem Grund größtenteils den Beobachtungen zuschreiben kann. Dr. \textsc{Olbers} selbst bezeichnet die Beobachtung vom 20. März als sehr zweifelhaft.

Bei diesem Vergleich berücksichtigte Dr. \textsc{Gauss} sowohl. Nutation und Aberration als auch die Parallaxe. Letztere ist besonders bei den ersten Beobachtungen, als der Abstand des Kome%
\SourceLeaf{253}%
ten von der Erde nur ein Viertel der Entfernung der Sonne betrug, nicht ganz unerheblich. Er bezog sie daher umso lieber ein, da sich der Rechenaufwand dadurch kaum erhöht.
Bei diesen Rechnungen verwendete er das Verfahren aus dem vortrefflichen Aufsatz, der im vorliegenden Heft auf S. 385 [S. 94 dieses Bandes] abgedruckt ist. Für die Konstanten $A$, $B$, $C$, $a$, $b$, $c$ fand er zur unmittelbaren Berechnung von Rektaszension und Deklination die folgenden Werte:
\[
A = 268^\circ\ 14'\ 21'', \quad B = 176^\circ\ 29'\ 48'', \quad C = 2^\circ\ 19'\ 58''
\]
\[
a = 87\ \ 19\ \ 38, \quad b = 56\ \ 56\ \ 7, \quad c = 33\ \ 12\ \ 3
\]

\Needspace{22\baselineskip}
\indent Die \emph{Konstanten} beziehen sich bereits auf den wahren Äquator und schließen somit die Nutation ein, die während einer so kurzen Zeit als unveränderlich anzusehen ist. Ist also $v$ die wahre Anomalie des Kometen, so erhält man die Koordinaten $x$, $y$, $z$ unmittelbar durch die folgenden äußerst einfachen Formeln:

\[
\begin{aligned}
x &= \frac{\alpha \sin(v + 240^\circ\ 11'\ 14'')}{\cos\frac{1}{2}v^2}, \\[5pt]
y &= \frac{\ObservedBeta \sin(v + 148^\circ\ 26'\ 41'')}{\cos\frac{1}{2}v^2}, \\[5pt]
z &= \frac{\gamma \sin(v + 334^\circ\ 16'\ 51'')}{\cos\frac{1}{2}v^2}
\end{aligned}
\]

\noindent und wobei

\[
\log\alpha = 0.029384, \quad \log \ObservedBeta = 9.953130, \quad \log \gamma = 9.768302
\]

\indent Die Größen $\alpha$, $\ObservedBeta$, $\gamma$ sind hier nämlich die Produkte aus dem Abstand im Perihel und dem Sinus von $a$, $b$, $c$. Die Winkel $240^\circ\ 11'\ 14''$ usw. sind gleich den Winkeln $A$, $B$, $C$ abzüglich der Entfernung des Knotens vom Perihel $=\,28^\circ\ 3'\ 7''$. — — —

\begin{flushright}
[\textsc{von Zach}.]
\end{flushright}

\bigskip
